\section{Arbitrarily light unbounded tails do not suffice at fixed temperature}
\label{sec:general-amplitude-obstruction}
The preceding Gaussian-mixture construction is not specific to Gaussian
amplitude tails. Even much lighter unbounded tails cannot, by themselves,
imply a general fixed-temperature raw-prediction oracle. The following
is still a generic function-prior example, not a counterexample for a
particular finite neural-network prior.

\begin{theorem}[One fixed prior for every fixed temperature]
\label{thm:general-amplitude-obstruction}
Let $G$ have any symmetric full-support law on $\mathbb R$ with
$\mathbb E|G|<\infty$. There is one centered prior $\mu$ on smooth
functions on $[0,1]$, independent of temperature, such that
\[
 \mathbb E_\mu\Psi(\|f\|_\infty)
       \le\mathbb E\Psi(|G|)
\]
for every nonnegative increasing function $\Psi$.
For every fixed $T>0$ there is a sequence of noiseless data laws,
with labels in $\{0,1\}$ and sample sizes $n_j\to\infty$, for which
the ordinary raw Gaussian-likelihood posterior mean satisfies
\[
 \mathbb P\{R(\bar f_{S_{n_j}})\ge c_Tj^2\}\longrightarrow1.
\]
The zero predictor has risk $1/2$ and a comparison law of entropy
at most $\log2$ under this one prior. In particular the conclusion applies to
$G=\operatorname{arsinh}(Z)$ with $Z$ standard Gaussian, whose moments
grow only logarithmically. The adverse law sequence is allowed to
depend on $T$; one fixed law or one law sequence defeating all
$T$ is not asserted.
\end{theorem}
\begin{proof}
For each integer $j\ge1$, set
$r_j=\mathbb P\{j\le G\le9j/8\}>0$.
Choose an integer
\[
 n_j\ge2^j+j^2\{1+\log(1/r_j)\},\qquad M_j=n_j^3,
\]
increasing with $j$, if necessary by increasing it further.
For each $k\ge1$ let $p_k=2^{-k}/1000$.
Choose disjoint open intervals $I_{j,k}\subset(0,1)$ and, inside
each, $M_j+1$ disjoint nonnegative smooth bump functions
$\zeta_{j,k,i}$ of maximum one. Let $x_{j,k,i}$ be points at which
the respective bumps equal one. For independent Bernoulli$(p_k)$
bits $B_1,\ldots,B_{M_j}$, define
\[
 \psi_{j,k,B}(x)=\frac1j\zeta_{j,k,0}(x)
                         +\sum_{i=1}^{M_j}B_i\zeta_{j,k,i}(x).
\]
The prior has an atom of mass $1/2$ at zero and, for each $(j,k)$,
mass $w_{j,k}=2^{-j-k-1}$ on $f=G\psi_{j,k,B}$, independently
drawing $G$ and the bits. These weights sum to one. Every individual
function is smooth and supported in one interval, and
$\|\psi_{j,k,B}\|_\infty\le1$. Symmetry centers the prior and
proves the asserted tail comparison.

Fix $T>0$ and choose $k$ such that
$p=p_k\le1/[1000(T+1)]$. For each $j$, put probability $1/2$ on
$x_{j,k,0}$ with label one, and probability $1/(2M_j)$ on each
$x_{j,k,i}$, $i\ge1$, with label zero. Every inactive prior branch
is identically zero at these inputs. Let $K$ count the common input
in a sample of size $n=n_j$. The event $K\ge n/3$ has probability
tending to one. We work on this event.

Within the active branch, integrating observed bits gives the
amplitude likelihood
\[
 \exp[-K(G/j-1)^2/T]L(G),\qquad
 L(G)=\prod_{i\,\mathrm{seen}}
       \{(1-p)+p\exp(-k_iG^2/T)\},
\]
where $k_i$ is the observed multiplicity of rare input $i$.
For every $G$, $(1-p)^n\le L(G)\le1$.
Its branch normalizer $D_j$ is therefore bounded below by
\[
 D_j\ge r_j\exp[-n/(64T)](1-p)^n
       \ge r_j\exp[-n/(64T)-2pn].
\]
If $G\le j/2$, the common-input squared error is at least $1/4$.
Consequently the active-branch posterior $Q_j$ satisfies
\[
 Q_j\{G\le j/2\}
 \le r_j^{-1}\exp[-n/(12T)+n/(64T)+2pn]
 \le r_j^{-1}e^{-n/(20T)}.
\]
For negative amplitudes the common-input error is at least one, so
\[
 \mathbb E_{Q_j}[|G|\mathbf1_{\{G<0\}}]
       \le(\mathbb E|G|)r_j^{-1}e^{-n/(4T)}.
\]
The last two numerical bounds follow from
$p\le1/[1000(T+1)]$. Since
$\log(1/r_j)\le n_j/j^2$, both right sides tend to zero.
It follows that $\mathbb E_{Q_j}G\ge j/4$ for all sufficiently
large $j$.

The combined inactive branches have likelihood $e^{-K/T}$.
The posterior odds of the active branch are at least
\[
 \frac{w_{j,k}}{1-w_{j,k}}\frac{D_j}{e^{-K/T}}
       \ge\frac{w_{j,k}}{1-w_{j,k}}r_j e^{n/(4T)}
       \longrightarrow\infty.
\]
Here $-\log w_{j,k}=O(j+k)$, $n_j\ge2^j$, and
$-\log r_j\le n_j/j^2$. Thus active posterior mass is at least
one half eventually. At an unobserved rare input, its bit retains
its independent Bernoulli$(p)$ law, even after conditioning on the
amplitude and all observed data. The full raw posterior mean there is
therefore at least $p j/8$. At least $M_j-n_j$ rare inputs are
unobserved, and their total population mass is at least $1/4$.
The resulting squared risk is at least $p^2j^2/256$, proving the
claim with $c_T=p_k^2/256$.

The zero predictor has exactly risk $1/2$. The zero function has prior probability at least $1/2$, giving the
stated comparison entropy. All posterior
means exist because the likelihood is bounded and
$\mathbb E|G|<\infty$; the displayed lower bounds remain valid if
some population risks are infinite.
\end{proof}

The theorem preserves any prescribed native tail control obtainable
from a symmetric unbounded amplitude law. For the asinh-Gaussian
choice it is therefore compatible with the logarithmic-temperature
positive oracle: it only rules out a generic conclusion at fixed
$T$. Its function prior uses a countable mixture of feature families
for the counterexample; it is not the fixed-architecture neural prior
used in that positive result. The quantitative divergence here is
measured in $j$, and its sample sizes can be much larger than in the
Gaussian-amplitude example.
